Vamos tentar, de maneira genérica, descobrir a função geradora para qualquer recurção de ordem 2.

Começando da seguinte recurção genérica,
\begin{align*}
  a_{n+2} &= \alpha a_{n+1} + \beta a_{n}, \\
  a_0 &= c, \\
  a_1 &= d,
\end{align*}

podemos convertê-la em uma função geradora da seguinte forma

\begin{calcproof}[.]
  \start{a_{n+2} = \alpha a_{n+1} + \beta a_{n}}
  \step[\calciff]{Conversão para notação de streams}{
    a = c \prec (\alpha a + (d - \alpha c) \prec (\beta a))
  }
  \step[\calciff]{Conversão para função geradora recursiva}{
    A(x) = c + x(\alpha A(x) + d - \alpha c + x(\beta A(x)))
  }
  \step[\calciff]{Distributiva}{
    A(x) = c + (\alpha A(x) + d - \alpha c)x + \beta A(x)x^2
  }
  \step[\calciff]{Soma dos dois lados}{
    A(x) - \alpha A(x)x - \beta A(x)x^2 = c + (d - \alpha c)x
  }
  \step[\calciff]{Distributiva}{
    A(x)(1 - \alpha x - \beta x^2) = c + (d - \alpha c)x
  }
  \step[\calciff]{Divisão dos dois lados}{
    A(x) = \dfrac{c + (d - \alpha c)x}{1 - \alpha x - \beta x^2}
  }
\end{calcproof}

Com isso, sabemos encontrar a função geradora de qualquer recurção de ordem 2. Agora, vamos nos encaminhar para o método das frações parciais. Primeiro, temos que encontrar as raízes $R_1$ e $R_2$ da polinômio de segundo grau do denominador da função geradora. Aplicando a fórmula de Bhaskara, encontramos
\[
  R_1, R_2 = \dfrac{\alpha \pm \sqrt{\alpha^2 + 4 \beta}}{-2 \beta}
\]

A princípio, esse resultando não parece muito interessante. No entanto, podemos notar que, quando comparamos as raízes com os autovalores encontrados pelo método da diagonalização, temos:

\begin{calcproof}
  \start{R_i = \dfrac{\alpha \pm \sqrt{\alpha^2 + 4 \beta}}{-2 \beta}}
  \step[\calciff]{Termo em evidência}{
    R_i = \dfrac{\alpha \pm \sqrt{\alpha^2 + 4 \beta}}{2 } \cdot \dfrac{1}{-\beta}
  }\step[\calciff]{Definição de $\lambda_i$}{
    R_i = \dfrac{\lambda_i}{-\beta}
  }
\end{calcproof}

Assim, temos sabemos como mapear todas as raízes do denominador da função geradora para um autovalor da matriz de transição do sistema dinâmico.

Esse mapeamento se torna ainda mais interessante quando percebemos que

\begin{calcproof}
  \start{R_1 R_2 = \dfrac{\alpha + \sqrt{\alpha^2 + 4 \beta}}{-2 \beta} \cdot \dfrac{\alpha - \sqrt{\alpha^2 + 4 \beta}}{-2 \beta}}
  \step[\calciff]{Produto notável}{
    R_1 R_2 = \dfrac{\alpha^2 - (\alpha^2 + 4 \beta)}{4\beta^2}
  }\step[\calciff]{Álgebra}{
    R_1 R_2 = \dfrac{-4 \beta}{4\beta^2}
  }\step[\calciff]{Álgebra}{
    R_1 R_2 = \dfrac{-1}{\beta}
  }
\end{calcproof}

Assim, concluímos que
\[
  R_i = \lambda_i R_1 R_2.
\]

Dessa forma,
\begin{gather*}
  R_1 = \dfrac{1}{\lambda_2}, \\
  R_2 = \dfrac{1}{\lambda_1},
\end{gather*}
ou seja, as raízes são o inverso dos autovalores!

Agora podemos começar o método das frações parciais.

\begin{calcproof}
  \start{A(x) = \dfrac{c + (d - \alpha c)x}{1 - \alpha x - \beta x^2}}
  \step[\calciff]{Fatoração do denominador}{
    A(x) = \dfrac{c + (d - \alpha c)x}{(1-\lambda_1 x)(1-\lambda_2 x)}
  }
  \step[\calciff]{Frações parciais}{
    A(x) = \dfrac{k_1}{1-\lambda_1 x} + \dfrac{k_2}{1-\lambda_2 x}
  }
\end{calcproof}

Note que na direita temos duas funções geradoras notáveis:
\[
  \dfrac{1}{1 - cx}
\]
é a sequência das potências de $c$. Portanto, encontramos a fórmula fechada dessa recursão!

\[
  a_n = k_1 \lambda_1^n + k_2 \lambda_2^n.
\]

No entanto, ainda não conhecemos os coeficientes, e, além disso, para ver melhor a relação com o método da diagonalização, vamos ter que continuar massageando o método das frações parciais:

\begin{calcproof}
  \start{A(x) = \dfrac{c + (d - \alpha c)x}{(1-\lambda_1 x)(1-\lambda_2 x)}}
  \step[\calciff]{Frações parciais}{
    \dfrac{c + (d - \alpha c)x}{(1-\lambda_1 x)(1-\lambda_2 x)} = \dfrac{k_1}{1-\lambda_1 x} + \dfrac{k_2}{1-\lambda_2 x}
  }\step[\calciff]{Multiplicação pelo denominador comum}{
    c + (d - \alpha c)x = k_1(1-\lambda_2x) + k_2(1-\lambda_1x)
  }\step[\calciff]{Distributiva}{
    c + (d - \alpha c)x = k_1-k_1\lambda_2x + k_2-k_2\lambda_1x
  }\step[\calciff]{Agrupamento dos termos}{
    (k_1+k_2) + (-k_1\lambda_2-k_2\lambda_1)x
  }
\end{calcproof}

Dessa forma, chegamos no seguinte sistema linear:
\begin{gather*}
  \begin{cases}
    c &= k_1 + k_2, \\
    d - \alpha c &= -k_2\lambda_1 -k_1\lambda_2.
  \end{cases}
\end{gather*}

Note que podemos transformar a segunda igualdade de forma que
\begin{calcproof}
  \start{d - \alpha c = -k_2\lambda_1 -k_1\lambda_2}
  \step[\calciff]{Equação 1}{
    d - \alpha (k_1 + k_2) = -k_2\lambda_1 -k_1\lambda_2
  }
  \step[\calciff]{Soma dos dois lados}{
    d = -k_2\lambda_1 -k_1\lambda_2 + \alpha (k_1 + k_2)
  }
  \step[\calciff]{Distributiva}{
    d = -k_2\lambda_1 -k_1\lambda_2 + \alpha k_1 + \alpha k_2
  }
  \step[\calciff]{Distributiva}{
    d = k_2(-\lambda_1 + \alpha) + k_1(-\lambda_2 + \alpha)
  }
  \step[\calciff]{$\lambda_1 + \lambda_2 = \alpha$}{
    d = k_2\lambda_2 + k_1\lambda_1
  }
\end{calcproof}

Dessa forma, o sistema se torna
\begin{gather*}
  \begin{cases}
    c &= k_1 + k_2, \\
    d &= k_2\lambda_2 + k_1\lambda_1,
  \end{cases}
\end{gather*}
que é idêntico ao sistema da diagonalização!
